\section{Implications for biofilm treatment}
\label{sec:implications}

Section~\ref{sec:intro} opened with a clinical question. Behavior-modifying interventions for biofilm infections, and in particular compounds that induce bacterial dispersal away from a biofilm into a more well-mixed, planktonic state, have been proposed as a way to reduce a bacterial population's capacity to evolve resistance under antibiotic pressure. The proposal does not require the intervention to kill cells directly; the hope is that disrupting biofilm spatial structure alone is enough to disarm the population, so that the surviving cells become more vulnerable to standard antibiotic treatment and to host immunity. The intuition behind that hope is that biofilm spatial structure helps the population evolve resistance, so disrupting that structure should reduce resistance evolution. Our findings from Sections~\ref{sec:doubleedge} through~\ref{sec:robustness} let us evaluate that intuition directly.

The answer depends on the antibiotic dose regime and on whether the biofilm's interior is metabolically active. At doses just above the minimum inhibitory concentration but below the crossover dose (our model places this near $\dEdge \approx 1.4$ at default parameters), the biofilm with an active core rescues less often than a well-mixed population of the same size (Section~\ref{sec:doubleedge}, Figure~\ref{fig:dosecurves}). Dispersing the biofilm at this dose moves the bacterial population from a configuration in which resistance evolution is suppressed by spatial structure to one in which it is more efficient. The intervention is counterproductive: the patient would face a population more likely to evolve resistance after the intervention than the one they had before.

At doses above the crossover ($\dEdge > 1.4$ in our model), the biofilm with an active core rescues more often than a well-mixed population. Dispersing the biofilm at this dose moves the population in the opposite direction, from the configuration that favors resistance evolution to the one that suppresses it. The intervention does what the clinical hope requires: it reduces the chance that the surviving bacterial population evolves resistance. If the biofilm's interior is metabolically dormant, however, a third pattern holds. A dormant biofilm rescues less often than a well-mixed population at every dose tested (Section~\ref{sec:robustness}). Dispersing a dormant biofilm always increases the chance of resistance evolution, regardless of dose. The clinical hope that disrupting biofilm structure reduces resistance never holds in this case.

Several open questions can be addressed within this modeling framework. Two were promised in earlier sections and not yet pursued, but the existing simulator has the data to address them. First, why does the dormant biofilm produce the same number of mutations as the well-mixed reference despite living substantially longer (Figure~\ref{fig:decomposition} Panel A)? The simulator records the full population trajectory for every run, and the answer almost certainly lies in the integrated cell-time of the wild-type population: a longer-lived population that is, on average, smaller can produce the same total mutation budget as a shorter-lived population that is, on average, larger. Working this out from the model's deterministic limit would convert the supply equality from an observed coincidence into a derived consequence. Second, why does an edge-born mutation in a biofilm establish less often than the same mutation in a well-mixed population (Figure~\ref{fig:decomposition} Panel B)? The lineage tracking already in place would let us isolate the contributions of the conveyor belt's continuous replacement of wild-type cells (which displaces nascent R lineages with refreshed competitors) and of the density-factor coupling between core size and edge births (which keeps the edge near saturation in a way the well-mixed reference does not). Neither analysis requires changes to the simulator.

A third question is addressable with a small modification to the simulator: instead of comparing a biofilm to a separate well-mixed population, we could simulate the dispersal intervention itself in-run, by running the biofilm until a chosen intervention time and then converting it to a well-mixed (or partially dispersed) population mid-simulation. The rescue probability of such a hybrid run would be a direct in-model prediction of what a dispersal-inducing intervention does over the course of a treatment, including the dynamics of the transition itself.

A fourth set of questions requires extending the model. The 1D chain with a sharp step-function antibiotic profile is a stylization. Natural next versions include a 3D geometry with a graded drug-penetration profile, an explicit host immune response that finishes off the surviving wild-type population (Section~\ref{sec:intro}'s clinical hope quietly assumed this would happen), a heterogeneous core in which $\bCore$ varies with distance from the surface, and stochastic phenotypic switching between active and dormant states. Each adds complexity but stays within the same Gillespie framework, and each would let us test which assumptions the qualitative rescue flip actually depends on.

The model also makes two assumptions about the bacterial genotypes that, if relaxed, could change the qualitative results. First, every run starts with $\Nzero$ wild-type cells and zero resistant cells; new R lineages arise only through mutation during the simulation itself. Second, R cells have the same birth and death rates as wild-type cells in the core, so resistance carries no fitness cost in the absence of antibiotic. Real bacterial populations almost always carry low-frequency R variants as standing genetic variation before antibiotic exposure, and resistance often does carry a fitness cost when the antibiotic is absent: in a homogeneous well-mixed population, selection drives the standing R frequency downward over time. Recent work in spatial population genetics (Fruet and colleagues, Bitbol and colleagues, and related) has shown that spatial structure can act as a refuge that preserves these deleterious variants at higher frequencies than a well-mixed population would, by isolating them from competition with the fitter wild type. In our framework, this storage effect would be strongest at low $\bCore$, where a largely dormant interior allows R lineages to persist without being actively selected against. If we allowed standing variation and a fitness cost for R outside the edge, we would expect the dormant biofilm to rescue more often than the well-mixed reference at low $\bCore$, precisely the regime where our current model says the dormant biofilm has no rescue advantage at any dose. The dormant-loses-everywhere result is therefore likely specific to our no-standing-variation, no-fitness-cost assumptions; a more realistic treatment of pre-treatment dynamics would probably reverse the ordering of dormant biofilm and well-mixed at low $\bCore$, with storage effects becoming the dominant axis promoting rescue rather than the mutation-supply axis we have studied here.

A separate biological assumption concerns the antibiotic's mode of action. The model treats the drug as purely bactericidal: it raises the death rate of wild-type cells in the edge ($\dEdge$ becomes the swept variable) but leaves their birth rate $\bEdge$ unchanged. Many real antibiotics are bacteriostatic (they slow or halt growth without killing cells directly), and many work via a mix of both modes. Under a purely bacteriostatic drug, wild-type cells in the edge would stop dividing rather than dying, so the edge would shrink only slowly through deaths at the baseline rate, the boundary would retreat much more slowly, and the supply of new edge-born mutations would drop because fewer edge births occur. The race from Section~\ref{sec:reachingedge} between boundary retreat and core drift would still play out, but the boundary speed would no longer scale with dose in the way our bactericidal model assumes, and the dose-dependent shifts in the crossover could change in both magnitude and direction. Whether the qualitative rescue flip survives in a bacteriostatic regime is an open question the current simulator could answer with a small parameter change: making $\bEdge$ dose-sensitive instead of $\dEdge$.

The clinical question from Section~\ref{sec:intro} does not have a single answer. The biofilm's spatial structure helps the population evolve resistance only at high antibiotic dose and with a metabolically active interior; in every other regime we tested, it hurts. A dispersal-inducing intervention is therefore context-dependent: useful in some clinical conditions, counterproductive in others, with the race mechanism in Section~\ref{sec:reachingedge} determining which. The broader takeaway is that the double-edged sword from Section~\ref{sec:doubleedge} captures a mechanistically grounded contingency: spatial structure in an evolving bacterial population under drug pressure is neither uniformly helpful nor uniformly harmful, and the direction of its effect on resistance evolution depends on a tractable race between mutation supply in the protected interior and mutation transport to the drug-exposed surface.
